\chapter{Modelo del mecanismo}
\label{chap:modelo}

Este capítulo desarrolla el modelo estático del mecanismo: relaciona el giro del tambor con el
ángulo de la barra, la longitud del cable y su tensión, y de ahí obtiene el par. Con ese modelo se
comparan tres perfiles de tambor y se establece el método con el que, en el
capítulo~\ref{chap:resultados}, se contrastan las medidas. Las mismas ecuaciones están
implementadas en \fichero{static/mecanismo.js}, de modo que el simulador y el visor de ensayos
calculan exactamente lo mismo que aquí se escribe.

\section{Descripción del mecanismo}

La figura~\ref{fig:mecanismo} muestra el esquema. Una barra rígida articulada en O lleva suspendida
una masa $m$ a una distancia $L_3+L_4$ de la articulación. El cable se ancla a la barra en el punto
A, a una distancia $L_3$ de O, y sube hasta el punto de salida D, desde donde se enrolla en la
caracola. El tambor gira solidario al eje de salida de la transmisión, con una reducción $N$ de
vueltas de motor por vuelta de tambor.

El punto D \emph{no} está sobre la vertical de O: está a una altura $L_2$ y desplazado $dx$ hacia el
lado contrario a la barra. Ese detalle, que en una primera versión del modelo se pasó por alto,
cambia la física y no sólo el dibujo, como se ve en la sección~\ref{sec:mod-tension}.

Como variable angular se usa $\alpha$, el ángulo en O entre la línea O$\to$D y la barra, que recorre
de 0° (barra plegada sobre esa línea) a 180° (barra en su prolongación). El ángulo respecto a la
horizontal es $\theta = \beta - \alpha$, siendo $\beta$ la inclinación de O$\to$D. Se prefiere
$\alpha$ porque a cada longitud de cable le corresponde un único $\alpha$, sin la ambigüedad que
tenía medir desde la horizontal.

\begin{figure}[H]
\centering
\begin{tikzpicture}[scale=1.7,>=Stealth]
  \def\al{100}
  \coordinate (O) at (0,0);
  \coordinate (D) at (-0.65,2.9);
  \coordinate (A) at ({0.95*cos(12)},{0.95*sin(12)});
  \coordinate (W) at ({2.05*cos(12)},{2.05*sin(12)});
  \fill[gray!25] (-0.35,-0.12) rectangle (0.35,0);
  \draw[gray!55,line width=1.2pt,dashed] (O) -- (D);
  \draw[dashed,gray] (O) -- (1.3,0);
  \draw[line width=2.4pt] (O) -- (W);
  \draw[blue!70!black,thick] (D) -- (A) node[pos=0.45,right] {$c$};
  \draw[violet,thick] (D) circle (0.2);
  \draw[violet] plot[domain=0:600,samples=90,smooth]
       ({-0.65+(0.04+0.00025*\x)*cos(\x)},{2.9+(0.04+0.00025*\x)*sin(\x)});
  \node[violet,right] at (-0.4,3.2) {caracola $r(\varphi)$};
  \node[below right=-1pt,font=\small] at (-0.52,2.70) {D};
  \fill (D) circle (1.2pt);
  \fill (O) circle (1.6pt) node[below=3pt] {O};
  \fill[blue!70!black] (A) circle (1.4pt) node[below right=-1pt] {A};
  \draw[->] (0.45,0) arc (0:12:0.45);
  \node[font=\small] at (0.62,0.07) {$\theta$};
  \draw[->,green!45!black] ({0.33*cos(\al+12)},{0.33*sin(\al+12)}) arc (\al+12:12:0.33);
  \node[green!45!black,font=\small] at (0.02,0.46) {$\alpha$};
  \draw (W) -- ++(0,-0.38);
  \fill[orange!80!black] ($(W)+(-0.15,-0.66)$) rectangle ($(W)+(0.15,-0.38)$);
  \draw[->,thick] ($(W)+(0,-0.7)$) -- ++(0,-0.4) node[right] {$mg$};
  \draw[->,thick,red!75!black] (A) -- ($(A)!0.3!(D)$) node[left=2pt] {$T$};
  \draw[<->] (-1.15,0) -- (-1.15,2.9) node[midway,left] {$L_2$};
  \draw[dotted] (-1.2,2.9) -- (-0.7,2.9);
  \draw[<->] (-0.65,-0.35) -- (0,-0.35) node[midway,below] {$dx$};
  \node[above=1pt] at ($(O)!0.55!(A)$) {$L_3$};
  \node[above left=0pt] at ($(A)!0.5!(W)$) {$L_4$};
\end{tikzpicture}
\caption{Esquema del mecanismo. El punto de salida del cable D está desplazado $dx$ respecto de la
vertical de la articulación O.}
\label{fig:mecanismo}
\end{figure}

El modelo parte de las siguientes hipótesis:
\begin{enumerate}[label=H\arabic*.]
  \item Régimen cuasiestático: a velocidad constante se desprecian las fuerzas de inercia.
  \item Cable inextensible y sin masa, y punto de salida D fijo.
  \item Se desprecia la masa de la barra frente a la masa suspendida.
  \item Las pérdidas por rozamiento no entran en el modelo estático: se tratan aparte en la
        sección~\ref{sec:mod-separacion}.
\end{enumerate}

\section{Longitud del cable}

Con $|OD| = \sqrt{L_2^2 + dx^2}$, la ley del coseno en el triángulo O--D--A da
\begin{equation}
  c(\alpha)^2 = |OD|^2 + L_3^2 - 2\,|OD|\,L_3\cos\alpha .
  \label{eq:cable}
\end{equation}
El coseno es monótono en $[0^\circ,180^\circ]$, de modo que la correspondencia entre $\alpha$ y $c$
es biunívoca: cada longitud de cable define un único ángulo, y al revés. Los extremos geométricos
son $c_{\min} = \big||OD|-L_3\big|$ con la barra plegada y $c_{\max} = |OD|+L_3$ con la barra
alineada. Despejando,
\begin{equation}
  \cos\alpha = \frac{|OD|^2 + L_3^2 - c^2}{2\,|OD|\,L_3}.
  \label{eq:angulo}
\end{equation}

\section{Tensión del cable}
\label{sec:mod-tension}

El momento del peso respecto a O es $mg(L_3+L_4)\cos\theta$. El brazo de la tensión es la distancia
de O a la recta D--A, que vale
\begin{equation}
  d = \frac{|OD|\,L_3\,\operatorname{sen}\alpha}{c}.
\end{equation}
Igualando momentos, $T\,d = mg(L_3+L_4)\cos\theta$, resulta
\begin{equation}
  \boxed{\;T = \frac{m\,g\,(L_3+L_4)\;c\,\cos\theta}{|OD|\;L_3\,\operatorname{sen}\alpha}\;}
  \qquad \theta = \beta - \alpha,\quad \beta = \arctan\frac{L_2}{-dx}.
  \label{eq:tension}
\end{equation}

Conviene detenerse en el caso particular $dx=0$, que fue el primer modelo del proyecto. Entonces
$\beta = 90^\circ$, $\theta = 90^\circ-\alpha$, $\cos\theta = \operatorname{sen}\alpha$, y los dos
factores se cancelan:
\begin{equation}
  T\big|_{dx=0} = \frac{m\,g\,(L_3+L_4)}{L_2\,L_3}\,c = K\,c ,
  \label{eq:tension0}
\end{equation}
es decir, \emph{la tensión sería proporcional a la longitud libre de cable}. Es un resultado
elegante, y por eso costó abandonarlo: con D desplazada la cancelación desaparece y la tensión deja
de ser proporcional al cable. Además, en $\alpha \to 180^\circ$ el brazo del cable se anula mientras
la pesa sigue haciendo momento, así que la tensión crece sin límite: ese punto es una singularidad
del mecanismo y no un punto de trabajo.

La tabla~\ref{tab:geometria} recoge los valores usados. No son cotas medidas, sino los parámetros
con los que trabaja la configuración del banco, ajustados sobre los propios ensayos
(sección~\ref{sec:res-ajuste}).

\begin{table}[H]
\centering
\caption{Parámetros del mecanismo empleados en el modelo \verificar{medir las cotas reales del
prototipo y rehacer el ajuste; de ellas dependen todas las cifras de este capítulo}.}
\label{tab:geometria}
\begin{tabular}{@{}lll@{}}
\toprule
Parámetro & Valor & Comentario \\
\midrule
$L_2$ & \SI{290}{\milli\metre} & altura de D sobre O \\
$dx$ & \SI{65}{\milli\metre} & desplazamiento horizontal de D \\
$|OD|$ & \SI{297.2}{\milli\metre} & línea O$\to$D, a \SI{12.6}{\degree} de la vertical \\
$L_3$ & \SI{95}{\milli\metre} & anclaje del cable en la barra \\
$L_4$ & \SI{110}{\milli\metre} & resto de barra hasta la masa \\
$m$ & \SI{2}{\kilo\gram} & masa suspendida \\
$\alpha_0$ & \SI{114}{\degree} & ángulo en B, donde arranca el ensayo \\
$c$ en B & \SI{346.9}{\milli\metre} & de~\eqref{eq:cable} \\
$c_{\min}$ -- $c_{\max}$ & \SIrange{202.2}{392.2}{\milli\metre} & topes geométricos \\
$r_0 \to r_1$ & \SIrange{45}{25}{\milli\metre} & radios de la caracola configurados \\
Barrido & \SI{0.3}{rev} de tambor & recorrido en que el radio pasa de $r_0$ a $r_1$ \\
$N$ & 5,5 & reducción motor:tambor \\
\bottomrule
\end{tabular}
\end{table}

\section{Par en el tambor}

El tambor paga cable según su radio instantáneo. Tomando como origen el punto B, donde arranca el
ensayo, y contando el giro $\varphi$ positivo en el sentido en que el ciclo avanza (de B hacia A):
\begin{equation}
  s(\varphi) = \int_0^{\varphi} r(u)\,\mathrm{d}u, \qquad
  c(\varphi) = c_B - s(\varphi), \qquad
  \varphi = \frac{2\pi\,p}{N},
  \label{eq:cable-tambor}
\end{equation}
donde $p$ es el avance del eje del motor en revoluciones. El par en el tambor y en el motor son
\begin{equation}
  \tau(\varphi) = T\big(c(\varphi)\big)\; r(\varphi), \qquad \tau_m = \frac{\tau}{N\,\eta},
  \label{eq:par}
\end{equation}
con $\eta$ el rendimiento de la transmisión. Para la caracola se modela un radio lineal en el giro
(espiral de Arquímedes) entre $r_0$ y $r_1$ a lo largo del barrido $\varphi_b$:
\begin{equation}
  r(\varphi) = r_0 + (r_1-r_0)\,\frac{\varphi}{\varphi_b}, \qquad 0 \le \varphi \le \varphi_b,
  \label{eq:espiral}
\end{equation}
y fuera de ese tramo el radio se queda en el extremo que corresponda.

\subsection{Cota inferior del par de pico}

El trabajo que el tambor entrega a lo largo de un recorrido es el que gana la masa en energía
potencial, y no depende del perfil:
\begin{equation}
  W = \int_0^{\varphi_{\text{tot}}} \tau\,\mathrm{d}\varphi = \int_{c_{\text{fin}}}^{c_{\text{ini}}} T(c)\,\mathrm{d}c .
  \label{eq:trabajo}
\end{equation}
Como la integral del par está fijada, su máximo no puede ser menor que su media:
\begin{equation}
  \max_{\varphi}\tau(\varphi) \;\geq\; \frac{W}{\varphi_{\text{tot}}},
  \label{eq:cota}
\end{equation}
y la igualdad sólo se alcanza si el par es constante. \emph{El perfil de par constante es, por
tanto, el que minimiza el par de pico para un recorrido y un giro de tambor dados.} Este resultado
no depende de la geometría: vale igual con $dx=0$ que sin él.

\subsection{Perfil ideal de par constante}

Imponiendo $\tau = T(c)\,r = \tau_0$ y usando $\mathrm{d}c/\mathrm{d}\varphi = -r$:
\begin{equation}
  r(\varphi) = \frac{\tau_0}{T\big(c(\varphi)\big)}, \qquad
  \frac{\mathrm{d}c}{\mathrm{d}\varphi} = -\frac{\tau_0}{T(c)},
  \qquad \tau_0 = \frac{W}{\varphi_{\text{tot}}} .
  \label{eq:ideal}
\end{equation}
Con $dx=0$, donde $T=Kc$, la ecuación se integra en forma cerrada y da
$c(\varphi)=\sqrt{c_0^2-2(\tau_0/K)\varphi}$, con el radio inversamente proporcional al cable. Con
$dx\neq0$ no hay forma cerrada y la integración se hace numéricamente, que es como se ha calculado
la curva de la figura~\ref{fig:perfiles}.

La regla práctica se mantiene en ambos casos: \emph{el radio debe ser grande donde la tensión es
pequeña y pequeño donde la tensión es grande}.

\subsection{Comparación de perfiles}

La comparación se hace sobre el recorrido que los ensayos hacen realmente con carga en el cable
(sección~\ref{sec:res-campania}): un avance del motor de \num{0.65} a \SI{1.80}{rev}, que son
\SI{0.209}{rev} de tambor y llevan el cable de \num{316.4} a \SI{276.6}{\milli\metre}. En ese tramo
el modelo da una tensión que \emph{baja} de \num{44.2} a \SI{35.5}{\newton}, porque el brazo del
cable mejora más deprisa de lo que crece el momento del peso, y un trabajo $W = \SI{1.58}{\joule}$.

\begin{table}[H]
\centering
\caption{Perfiles de tambor sobre el recorrido con carga ($\varphi_{\text{tot}}=\SI{0.209}{rev}$ de
tambor, $W=\SI{1.58}{\joule}$).}
\label{tab:perfiles}
\begin{tabular}{@{}lcccc@{}}
\toprule
Perfil & $r$ inicial & $r$ final & $\tau$ mín. -- máx. & Rizado \\
       & (mm) & (mm) & (N·m) & $\frac{\max-\min}{\text{media}}$ \\
\midrule
Cilindro equivalente & 30,3 & 30,3 & 1,06 -- 1,35 & \SI{24}{\percent} \\
Caracola configurada & 41,2 & 27,9 & 0,88 -- 1,65 & \SI{65}{\percent} \\
Ideal (par constante) & 27,2 & 33,8 & 1,20 -- 1,20 & \SI{0}{\percent} \\
\bottomrule
\end{tabular}
\end{table}

El resultado es llamativo y conviene subrayarlo: con esta geometría, y en este tramo, \emph{la
caracola tal y como está configurada empeora el reparto del par} respecto a un cilindro que
recorriese lo mismo (rizado del \SI{65}{\percent} frente al \SI{24}{\percent}). La razón es que la
tensión baja a lo largo del recorrido, de modo que para compensarla el radio tendría que
\emph{crecer} un \SI{24}{\percent}, y el perfil configurado lo hace decrecer un \SI{32}{\percent}.
Es decir: la caracola está descrita —o montada— al revés respecto de lo que pide la compensación.

\begin{figure}[H]
\centering
\includegraphics[width=\textwidth]{n_perfiles.pdf}
\caption{Izquierda: par medido (normalizado) frente a las dos orientaciones posibles del perfil.
Derecha: par en el tambor para el cilindro equivalente, la caracola configurada y el perfil ideal.}
\label{fig:perfiles}
\end{figure}

\section{Método experimental: separación del par y radio efectivo}
\label{sec:mod-separacion}

El accionamiento entrega el par en unidades internas de escala desconocida, y el par medido incluye
las pérdidas por rozamiento. Para comparar con el modelo sin depender de esa escala se recorre el
mismo punto en los dos sentidos, a la misma velocidad. Las pérdidas se oponen siempre al
movimiento, mientras que el par gravitatorio conserva el signo:
\begin{equation}
  |\tau_{\uparrow}(p)| = \tau_g(p) + \tau_f(p), \qquad
  |\tau_{\downarrow}(p)| = \tau_g(p) - \tau_f(p),
\end{equation}
\begin{equation}
  \tau_g = \frac{|\tau_{\uparrow}| + |\tau_{\downarrow}|}{2}, \qquad
  \tau_f = \frac{|\tau_{\uparrow}| - |\tau_{\downarrow}|}{2}.
  \label{eq:separacion}
\end{equation}
El mismo promedio aplicado a la señal de la célula elimina en primera aproximación el retardo de su
filtro, que desplaza en sentidos opuestos las curvas de ida y de vuelta.

Con $\tau_g$ ya limpio de pérdidas, la ecuación~\eqref{eq:par} permite estimar el \emph{radio
efectivo}
\begin{equation}
  r_{\text{ef}}(p) \;\propto\; \frac{\tau_g(p)}{T(p)},
  \label{eq:radio-efectivo}
\end{equation}
donde $T$ se toma del modelo, que no necesita la célula de carga: sólo la geometría y la masa. En un
tambor cilíndrico $r_{\text{ef}}$ sería constante; su forma es, por tanto, una medida directa del
perfil del tambor.

\subsection{Por qué no basta con medir la fuerza}
\label{sec:mod-ambiguedad}

Podría pensarse que basta con comparar la señal de la célula con la tensión del modelo. No es así, y
esto condiciona todo el capítulo de resultados. Como la célula no está calibrada, la comparación se
hace ajustando una recta $T \approx a\,(\text{cuentas}) + b$, y el coeficiente de determinación de
ese ajuste es el cuadrado del coeficiente de correlación:
\begin{equation}
  R^2 = \rho^2 .
\end{equation}
El valor de $R^2$ \emph{no cambia si se invierte el signo de $a$}. Una geometría que prediga una
tensión creciente y otra que la prediga decreciente pueden dar exactamente el mismo $R^2$ frente a
los mismos datos, y sólo difieren en el signo de la sensibilidad de la célula, que es precisamente
lo que no se conoce. La medida de fuerza, por sí sola, no puede decidir en qué sentido trabaja el
mecanismo.

El par del accionamiento sí puede: su signo y su escala son conocidos y no dependen de ninguna
calibración pendiente. Por eso el criterio que se usa en el capítulo~\ref{chap:resultados} es
comparar $\tau_g$ con el par del modelo, y no sólo la fuerza.

\section{Limitaciones del modelo}

\begin{itemize}
  \item Los parámetros de la tabla~\ref{tab:geometria} no son cotas medidas, sino un ajuste sobre
        los ensayos; el propio capítulo~\ref{chap:resultados} muestra que no son consistentes con el
        par medido.
  \item El punto de salida del cable se desplaza con el radio de enrollamiento, lo que introduce un
        error de algunos milímetros en $L_2$ y $dx$.
  \item Se desprecian la elasticidad del cable y la masa de la barra.
  \item El rendimiento y las pérdidas se suponen iguales en los dos sentidos, en valor absoluto,
        para poder aplicar~\eqref{eq:separacion}.
  \item El modelo es estático: no describe el tramo inicial en que el cable está destensado, ni los
        transitorios de arranque e inversión.
\end{itemize}
